\section{Complete Training Pipeline Example}

\begin{lstlisting}[style=python, caption=train\_pipeline.py]
import mlflow
import pandas as pd
from sklearn.model_selection import train_test_split
from sklearn.ensemble import RandomForestClassifier
from sklearn.metrics import accuracy_score, classification_report
import joblib

def load_data(path):
    """Load and prepare data"""
    df = pd.read_csv(path)
    X = df.drop('target', axis=1)
    y = df['target']
    return train_test_split(X, y, test_size=0.2, random_state=42)

def train_model(X_train, y_train, params):
    """Train RandomForest model"""
    model = RandomForestClassifier(**params)
    model.fit(X_train, y_train)
    return model

def evaluate_model(model, X_test, y_test):
    """Evaluate model performance"""
    y_pred = model.predict(X_test)
    accuracy = accuracy_score(y_test, y_pred)
    report = classification_report(y_test, y_pred)
    return accuracy, report

def main():
    # MLflow experiment tracking
    mlflow.set_experiment("model_training")

    with mlflow.start_run():
        # Load data
        X_train, X_test, y_train, y_test = load_data('data/processed/data.csv')

        # Model parameters
        params = {
            'n_estimators': 100,
            'max_depth': 10,
            'random_state': 42
        }
        mlflow.log_params(params)

        # Train model
        model = train_model(X_train, y_train, params)

        # Evaluate model
        accuracy, report = evaluate_model(model, X_test, y_test)
        mlflow.log_metric("accuracy", accuracy)

        # Save model
        mlflow.sklearn.log_model(model, "model")
        joblib.dump(model, 'models/model.pkl')

        print(f"Accuracy: {accuracy:.4f}")
        print(report)

if __name__ == "__main__":
    main()
\end{lstlisting}

\section{REST API for Model Serving}

\begin{lstlisting}[style=python, caption=serve.py]
from fastapi import FastAPI, HTTPException
from pydantic import BaseModel
import joblib
import numpy as np
import logging

# Setup logging
logging.basicConfig(level=logging.INFO)
logger = logging.getLogger(__name__)

# Load model
model = joblib.load('models/model.pkl')

# Create FastAPI app
app = FastAPI(title="ML Model API", version="1.0")

class PredictionRequest(BaseModel):
    features: list[float]

class PredictionResponse(BaseModel):
    prediction: int
    probability: float

@app.get("/health")
async def health_check():
    """Health check endpoint"""
    return {"status": "healthy"}

@app.post("/predict", response_model=PredictionResponse)
async def predict(request: PredictionRequest):
    """Make prediction"""
    try:
        # Prepare input
        X = np.array(request.features).reshape(1, -1)

        # Make prediction
        prediction = model.predict(X)[0]
        probability = model.predict_proba(X)[0].max()

        logger.info(f"Prediction: {prediction}, Probability: {probability:.4f}")

        return PredictionResponse(
            prediction=int(prediction),
            probability=float(probability)
        )
    except Exception as e:
        logger.error(f"Prediction error: {str(e)}")
        raise HTTPException(status_code=500, detail=str(e))

if __name__ == "__main__":
    import uvicorn
    uvicorn.run(app, host="0.0.0.0", port=8000)
\end{lstlisting}

\section{Data Validation with Great Expectations}

\begin{lstlisting}[style=python, caption=validate\_data.py]
import great_expectations as ge
import pandas as pd

def create_expectation_suite(df):
    """Create data validation suite"""
    # Convert to GE DataFrame
    ge_df = ge.from_pandas(df)

    # Define expectations
    ge_df.expect_table_row_count_to_be_between(min_value=100, max_value=1000000)
    ge_df.expect_column_values_to_not_be_null('user_id')
    ge_df.expect_column_values_to_be_unique('user_id')
    ge_df.expect_column_values_to_be_between('age', min_value=0, max_value=120)
    ge_df.expect_column_values_to_be_in_set('country', ['US', 'UK', 'CA', 'AU'])
    ge_df.expect_column_mean_to_be_between('revenue', min_value=0, max_value=10000)

    # Save expectation suite
    ge_df.save_expectation_suite('data_expectations.json')

    return ge_df

def validate_data(df, expectation_suite_path):
    """Validate data against expectations"""
    ge_df = ge.from_pandas(df)

    # Load expectations
    ge_df.read_json(expectation_suite_path)

    # Validate
    results = ge_df.validate()

    return results

if __name__ == "__main__":
    # Load data
    df = pd.read_csv('data/raw/data.csv')

    # Create and validate
    ge_df = create_expectation_suite(df)
    results = validate_data(df, 'data_expectations.json')

    print(f"Validation success: {results.success}")
    print(f"Number of expectations: {len(results.results)}")
\end{lstlisting}

\section{Airflow DAG Template}

\begin{lstlisting}[style=python, caption=ml\_dag.py]
from airflow import DAG
from airflow.operators.python import PythonOperator
from airflow.operators.bash import BashOperator
from datetime import datetime, timedelta
import sys

sys.path.insert(0, '/opt/airflow/dags/src')

from data_processing import extract_data, transform_data
from model_training import train_model
from model_evaluation import evaluate_model

default_args = {
    'owner': 'ml_team',
    'depends_on_past': False,
    'start_date': datetime(2024, 1, 1),
    'email': ['ml-team@example.com'],
    'email_on_failure': True,
    'email_on_retry': False,
    'retries': 3,
    'retry_delay': timedelta(minutes=5),
}

with DAG(
    'ml_training_pipeline',
    default_args=default_args,
    description='Complete ML training pipeline',
    schedule_interval='@daily',
    catchup=False,
    tags=['ml', 'training'],
) as dag:

    extract_task = PythonOperator(
        task_id='extract_data',
        python_callable=extract_data,
    )

    transform_task = PythonOperator(
        task_id='transform_data',
        python_callable=transform_data,
    )

    validate_task = BashOperator(
        task_id='validate_data',
        bash_command='python /opt/airflow/dags/src/validate_data.py',
    )

    train_task = PythonOperator(
        task_id='train_model',
        python_callable=train_model,
    )

    evaluate_task = PythonOperator(
        task_id='evaluate_model',
        python_callable=evaluate_model,
    )

    deploy_task = BashOperator(
        task_id='deploy_model',
        bash_command='python /opt/airflow/dags/src/deploy_model.py',
    )

    extract_task >> transform_task >> validate_task >> train_task >> evaluate_task >> deploy_task
\end{lstlisting}

\section{Model Monitoring Script}

\begin{lstlisting}[style=python, caption=monitor.py]
import pandas as pd
import numpy as np
from scipy import stats
from evidently.report import Report
from evidently.metric_preset import DataDriftPreset, DataQualityPreset

def detect_data_drift(reference_data, current_data, threshold=0.05):
    """Detect data drift using KS test"""
    drift_detected = {}

    for column in reference_data.columns:
        if pd.api.types.is_numeric_dtype(reference_data[column]):
            statistic, pvalue = stats.ks_2samp(
                reference_data[column].dropna(),
                current_data[column].dropna()
            )
            drift_detected[column] = {
                'drift': pvalue < threshold,
                'pvalue': pvalue,
                'statistic': statistic
            }

    return drift_detected

def generate_monitoring_report(reference_data, current_data):
    """Generate comprehensive monitoring report"""
    report = Report(metrics=[
        DataDriftPreset(),
        DataQualityPreset(),
    ])

    report.run(reference_data=reference_data, current_data=current_data)
    report.save_html('monitoring_report.html')

    return report

if __name__ == "__main__":
    # Load reference and current data
    reference = pd.read_csv('data/reference/data.csv')
    current = pd.read_csv('data/current/data.csv')

    # Detect drift
    drift_results = detect_data_drift(reference, current)

    # Print results
    for column, result in drift_results.items():
        if result['drift']:
            print(f"⚠ Drift detected in {column}: p-value={result['pvalue']:.4f}")

    # Generate detailed report
    report = generate_monitoring_report(reference, current)
\end{lstlisting}
